\documentclass[11pt]{article}

\usepackage[utf8]{inputenc}
\usepackage[T1]{fontenc}
\usepackage{lmodern}
\usepackage[margin=1in]{geometry}
\usepackage{amsmath,amssymb,bm}
\usepackage{booktabs}
\usepackage{array}
\usepackage[dvipsnames]{xcolor}
\usepackage{graphicx}
\usepackage{siunitx}
\usepackage{pgfplots}
\usepackage{pgfplotstable}
\usepgfplotslibrary{groupplots}
\pgfplotsset{compat=1.18}
\usepackage{placeins}
\usepackage{hyperref}
\usepackage{microtype}
\setlength{\emergencystretch}{3em}
\clubpenalty=10000
\widowpenalty=10000
\displaywidowpenalty=10000
\renewcommand{\topfraction}{0.95}
\renewcommand{\bottomfraction}{0.9}
\renewcommand{\textfraction}{0.05}
\renewcommand{\floatpagefraction}{0.8}
\setcounter{topnumber}{3}
\setcounter{totalnumber}{4}
\usepackage[backend=biber,style=numeric,sorting=none]{biblatex}
\addbibresource{References.bib}

\hypersetup{
    colorlinks=true,
    linkcolor=blue,
    citecolor=blue,
    urlcolor=blue
}

\newcommand{\dd}{\,\mathrm{d}}
\newcommand{\norm}[1]{\left\lVert #1 \right\rVert}
\newcommand{\abs}[1]{\left\lvert #1 \right\rvert}
\newcommand{\calI}{\mathcal{I}}
\newcommand{\calP}{\mathcal{P}}
\newcommand{\calG}{\mathcal{G}}
\newcommand{\eps}{\varepsilon}
\newcommand{\calK}{\mathcal{K}}

\newcommand{\DataDir}{data}
\newcommand{\GeneratedDir}{generated}


\pgfplotsset{
    transfer convergence axis/.style={
        xmode=log,
        ymode=log,
        x dir=reverse,
        grid=both,
        minor grid style={gray!15},
        major grid style={gray!30},
        xlabel={Effective spacing $h$},
        tick label style={font=\scriptsize},
        label style={font=\small},
        title style={font=\small},
        legend style={font=\scriptsize,draw=none},
        line width=0.8pt,
        mark size=2.0pt,
    },
    transfer bar axis/.style={
        grid=major,
        tick label style={font=\scriptsize},
        label style={font=\small},
        legend style={font=\scriptsize,draw=none},
    },
    scheme NN1/.style={black,solid,mark=*},
    scheme NN2/.style={NavyBlue,solid,mark=square*},
    scheme NN3/.style={ForestGreen,solid,mark=triangle*},
    scheme BAR2/.style={BrickRed,dashed,mark=o},
    scheme BAR3/.style={Orange,dashed,mark=square},
    scheme GP2/.style={Purple,densely dotted,mark=diamond*},
    scheme GP3/.style={Magenta,densely dotted,mark=triangle},
    reference slope/.style={gray!70,thin,dash pattern=on 3pt off 2pt,no marks},
}

\sisetup{
    scientific-notation=true,
    round-mode=places,
    round-precision=3,
}




% Keyed access to the existing control data. No numerical data are duplicated.
\pgfplotstableread[col sep=comma]{generated/reassessed_rans_controls_compact.csv}\ControlData
\pgfplotstableforeachcolumnelement{closure}\of\ControlData\as\ThisClosure{%
  \edef\SavedClosure{\ThisClosure}%
  \pgfplotstablegetelem{\pgfplotstablerow}{control}\of\ControlData\edef\ThisControl{\pgfplotsretval}%
  \pgfplotstablegetelem{\pgfplotstablerow}{quantity}\of\ControlData\edef\ThisQuantity{\pgfplotsretval}%
  \pgfplotstablegetelem{\pgfplotstablerow}{coeff}\of\ControlData\edef\ThisCoeff{\pgfplotsretval}%
  \pgfplotstablegetelem{\pgfplotstablerow}{production_rms}\of\ControlData
  \expandafter\xdef\csname Ctrl@\SavedClosure @\ThisControl @\ThisQuantity @\ThisCoeff @base\endcsname{\pgfplotsretval}%
  \pgfplotstablegetelem{\pgfplotstablerow}{change_rms}\of\ControlData
  \expandafter\xdef\csname Ctrl@\SavedClosure @\ThisControl @\ThisQuantity @\ThisCoeff @change\endcsname{\pgfplotsretval}%
  \pgfplotstablegetelem{\pgfplotstablerow}{change_over_production}\of\ControlData
  \expandafter\xdef\csname Ctrl@\SavedClosure @\ThisControl @\ThisQuantity @\ThisCoeff @ratio\endcsname{\pgfplotsretval}%
}
\newcommand{\Control}[5]{\csname Ctrl@#1@#2@#3@#4@#5\endcsname}
\newcommand{\ControlPercent}[4]{%
  \pgfmathparse{100*\Control{#1}{#2}{#3}{#4}{ratio}}%
  \pgfmathprintnumber[fixed,precision=1,fixed zerofill]{\pgfmathresult}}
\newcommand{\MaxControlPercent}[3]{%
  \pgfmathparse{100*max(\Control{#1}{#2}{NN2}{#3}{ratio},\Control{#1}{#2}{GP2}{#3}{ratio})}%
  \pgfmathprintnumber[fixed,precision=1,fixed zerofill]{\pgfmathresult}}
\newcommand{\ControlRow}[3]{#1 & #3 &
  \MaxControlPercent{#1}{#2}{CL} & \MaxControlPercent{#1}{#2}{CD} &
  \ControlPercent{#1}{#2}{P}{CL} & \ControlPercent{#1}{#2}{P}{CD}\\}
\pgfplotstableread[col sep=comma]{RANS_PostProcessing/rans_control_inner.csv}\InnerData
\pgfplotstableforeachcolumnelement{model}\of\InnerData\as\ThisModel{%
  \edef\SavedModel{\ThisModel}%
  \pgfplotstablegetelem{\pgfplotstablerow}{setting}\of\InnerData\edef\ThisSetting{\pgfplotsretval}%
  \pgfplotstablegetelem{\pgfplotstablerow}{method}\of\InnerData\edef\ThisMethod{\pgfplotsretval}%
  \pgfplotstablegetelem{\pgfplotstablerow}{steps}\of\InnerData
  \expandafter\xdef\csname In@\SavedModel @\ThisSetting @\ThisMethod @steps\endcsname{\pgfplotsretval}%
  \pgfplotstablegetelem{\pgfplotstablerow}{steps_at_cap}\of\InnerData
  \expandafter\xdef\csname In@\SavedModel @\ThisSetting @\ThisMethod @cap\endcsname{\pgfplotsretval}%
  \pgfplotstablegetelem{\pgfplotstablerow}{log10_rms_rho_end_median}\of\InnerData
  \expandafter\xdef\csname In@\SavedModel @\ThisSetting @\ThisMethod @res\endcsname{\pgfplotsretval}%
}
\newcommand{\Inner}[4]{\csname In@#1@#2@#3@#4\endcsname}
\newcommand{\InnerRow}[4]{#1 & #3 & #4 & \Inner{#1}{#2}{GP2}{steps} &
 \Inner{#1}{#2}{GP2}{cap}/\Inner{#1}{#2}{NN2}{cap} &
 \Inner{#1}{#2}{GP2}{res}/\Inner{#1}{#2}{NN2}{res}\\}
\pgfplotstableset{publication table/.style={
 col sep=comma,
 every head row/.style={before row=\toprule,after row=\midrule},
 every last row/.style={after row=\bottomrule}}}

\title{Mesh-to-Mesh Transfer Errors in Multistep Unsteady Flow Simulations}
% Author-supplied: replace the placeholders with the author list and affiliations.
\author{\textit{[author names]}\\ \textit{[affiliations]}}
\date{}
\begin{document}
\maketitle

\begin{abstract}
Mesh replacement changes both the spatial discretization and the solution histories used to advance an unsteady flow. This study relates mesh-to-mesh transfer errors to the subsequent state and aerodynamic response in second-order backward differentiation. Nearest-neighbor Taylor interpolation, barycentric interpolation, and continuous Galerkin projection are compared using smooth and discontinuous fields, repeated Euler remeshing, and stationary-airfoil calculations. Smooth tests recover the intended first- to third-order behavior. In the repeated-transfer tests, enforcing the global density integral to roundoff does not prevent substantial transition broadening, a result reproduced across three mesh-deformation realizations. Controlled airfoil restarts show that an error confined to the older history produces a finite next-step response even when the latest history is exact. Boundary enforcement also changes the transferred momentum entering the receiving calculation. With two Reynolds-averaged Navier--Stokes closures, Galerkin projection reduces turbulence-integral defects by factors of approximately 37--247, but the smaller differences in aerodynamic loads remain sensitive to time-step refinement and, for one closure, incomplete inner convergence. Transfer accuracy, admissibility, integral preservation, and multistep consistency must therefore be assessed separately; none alone establishes aerodynamic accuracy after remeshing.
\end{abstract}

\section{Introduction}
\label{sec:introduction}

Mesh adaptation can reduce the cost of unsteady flow calculations by concentrating resolution near evolving structures. When the mesh is replaced, however, the calculation must continue from a state represented on a different spatial discretization. A multistep integrator also requires previous solution levels. The accuracy of the continuation therefore depends not only on the new mesh and interpolated field, but on the interaction between the transferred histories and the receiving discrete equations. Time-accurate anisotropic adaptation and space--time adaptation for implicit Reynolds-averaged Navier--Stokes (RANS) calculations expose this coupling~\cite{AlauzetLoseilleOlivier2018,SauvageAlauzetDervieux2024}.

Accurate interpolation does not necessarily preserve a discrete balance, and conservation of a domain integral does not determine the spatial distribution of the transferred quantity. Common-refinement Galerkin projection preserves the consistent integral and minimizes the distance to the source field in the target finite-element space~\cite{FarrellMaddison2011,JiaoHeath2004}. Additional constraints are needed to preserve other discrete properties~\cite{MaddisonHiester2017}. Balance-preserving transfer in shallow-water models provides a related example: a small interpolation error can still disturb a steady numerical state~\cite{MaddisonCotterFarrell2011,MaddisonCotterFarrell2013}. For aerodynamic calculations, the relevant question is how such a disturbance affects the subsequent solution and loads.

This work examines that question for vertex-centered finite-volume calculations advanced with the second-order backward differentiation formula (BDF2). It compares nearest-neighbor Taylor interpolation, barycentric interpolation on a simplicial representation of the mesh, and continuous Galerkin projection, with nominal orders from one to three. Optional global correction separates preservation of an integral from the underlying pointwise interpolation. Particular attention is given to the distinction between the complete projected field and the nodal state accepted by the flow solver after restriction, limiting, and boundary enforcement.

The contribution is a controlled assessment of transfer-induced imbalance and multistep restart response, rather than a new interpolation construction. Static tests establish approximation and boundedness properties; repeated Euler remeshing tests their accumulation; and independent airfoil meshes permit the two BDF2 history contributions to be varied separately. Fixed-body RANS calculations then add turbulence transport and viscous loads. Prescribed mesh replacements keep the method comparisons independent of solution-driven adaptation. The physical domain remains fixed, and no continuous mesh motion is considered.

\section{Solution-transfer operators}
\label{sec:numerical-framework}

\subsection{Representation and derivative reconstruction}
\label{sec:derivative-reconstruction}

Source and target meshes, $A$ and $B$, cover the same discretized domain $\Omega$. The nodal state on mesh $M$ is $U_M$, and its finite-volume mass matrix is $M_M=\operatorname{diag}(w_i^M)\otimes I$, where the control-volume weights $w_i^M$ are positive. For one scalar component, the solver integral is
\begin{equation}
 Q_M(\phi_M)=\sum_i w_i^M\phi_{M,i}
 \label{eq:fv-discrete-integral}
\end{equation}
All methods return values at the flow-solver vertices. Mean-flow density, momentum, and total energy are transferred directly in conserved-variable form. Nodal interpolation is not control-volume averaging, even when Eq.~\eqref{eq:fv-discrete-integral} is used to assess its conservation.

Higher-order transfer requires derivatives recovered from neighboring values. A weighted least-squares fit uses
\begin{equation}
 \phi_j-\phi_i\simeq g_i\cdot r_{ij}+\tfrac12r_{ij}^{T}H_i r_{ij},
 \qquad r_{ij}=x_j-x_i
 \label{eq:wls-polynomial}
\end{equation}
A degree-one fit omits the Hessian. With smooth data, a full-rank scaled stencil, and controlled mesh geometry, the linear fit gives an $O(h)$ gradient error; the quadratic fit gives $O(h^2)$ gradient and $O(h)$ Hessian errors. Anisotropy, weighting, and stencil conditioning can degrade these estimates~\cite{BarthFrederickson1990,Barth1993,Mavriplis2003,DiskinThomas2008}. Third-order variants use the quadratic reconstruction, and their achieved order is verified numerically.

\subsection{Nearest-neighbor and barycentric interpolation}
\label{sec:point-transfer}
\label{sec:barycentric-transfer}

Let $n(j)$ be the source vertex closest to target point $x_j^B$, and $r_j=x_j^B-x_{n(j)}^A$. Nearest-neighbor Taylor transfer gives
\begin{align}
 (\calI^{\mathrm{NN1}}_{AB}\phi_A)_j&=\phi_{A,n(j)}\label{eq:nn1}\\
 (\calI^{\mathrm{NN2}}_{AB}\phi_A)_j&=\phi_{A,n(j)}+g_{n(j)}\cdot r_j\label{eq:nn2}\\
 (\calI^{\mathrm{NN3}}_{AB}\phi_A)_j&=\phi_{A,n(j)}+g_{n(j)}\cdot r_j
              +\tfrac12r_j^TH_{n(j)}r_j\label{eq:nn3}
\end{align}
The variants respectively copy the donor value, include its local slope, and include its curvature. For $|r_j|=O(h)$ and sufficiently accurate derivatives, their smooth-field errors are $O(h)$, $O(h^2)$, and $O(h^3)$. The donor changes across source-node Voronoi boundaries, so the transferred value need not vary continuously with target position.

Barycentric interpolation instead uses the source simplex containing the target point. A conforming triangulation or tetrahedralization of the original mesh provides the search geometry, without adding control-volume vertices. If the simplex vertices are $x_{i_r}^A$, the barycentric coordinates satisfy
\begin{equation}
 x_j^B=\sum_{r=1}^{d+1}\lambda_r x_{i_r}^A,
 \qquad \sum_r\lambda_r=1,
 \qquad \lambda_r\geq0
 \label{eq:barycentric-coordinates}
\end{equation}
The same geometric weights define the second-order interpolant and its gradient-corrected extension:
\begin{align}
 (\calI^{\mathrm{BAR2}}_{AB}\phi_A)_j
 &=\sum_r\lambda_r\phi_{A,i_r}\label{eq:bar2}\\
 (\calI^{\mathrm{BAR3}}_{AB}\phi_A)_j
 &=\sum_r\lambda_r\left[\phi_{A,i_r}
       +\tfrac12g_{i_r}\cdot(x_j^B-x_{i_r}^A)\right]\label{eq:bar3}
\end{align}
BAR2 reproduces affine fields and is bounded by the containing-simplex values. For a quadratic field with exact vertex gradients, the half-gradient term cancels its interpolation remainder; BAR3 is consequently third-order consistent. Shared values and gradients give continuity across conforming simplex interfaces. Neither method imposes a source-to-target integral constraint.

Derivative limiting changes these properties independently of nominal order. Multiplying the Taylor corrections by $0\leq\Phi_i\leq1$ reduces NN2 or NN3 to NN1 when $\Phi_i=0$, but reduces BAR3 to BAR2. The smooth verification is unlimited; discontinuous tests also use Barth--Jespersen limiting~\cite{BarthJespersen1989}. A limiter derived from linear increments does not automatically bound a quadratic reconstruction at every evaluation point. Extrema must therefore be checked for the actual returned field.

\subsection{Galerkin projection and cross-mesh integration}
\label{sec:galerkin-projection}
\label{sec:simplicial-representation}

Continuous Galerkin projection determines a coupled target field rather than evaluating target vertices independently. For a continuous degree-$p$ Lagrange space $V_B^p$, it satisfies
\begin{equation}
 (\phi_B^h-\phi_A^h,v_B)_{L_2(\Omega)}=0
 \quad\text{for every }v_B\in V_B^p
 \label{eq:galerkin-orthogonality}
\end{equation}
This is the best $L_2$ approximation of the source representation $\phi_A^h$ in $V_B^p$~\cite{FarrellMaddison2011,JiaoHeath2004}. Writing $\phi_B^h=\sum_jv_jN_j^B$ gives
\begin{equation}
 H_Bv=b,\qquad
 (H_B)_{ij}=\int_\Omega N_i^BN_j^B\dd x,
 \qquad b_i=\int_\Omega N_i^B\phi_A^h\dd x
 \label{eq:galerkin-system}
\end{equation}
The sparse positive-definite matrix $H_B$ couples basis functions sharing target elements. It differs from the diagonal finite-volume matrix $M_B$ used for time integration.

A target element can cross several source elements. Although a target basis function has one polynomial expression there, the source field changes expression at each source-element boundary. The cross-mesh load must therefore be integrated as
\begin{equation}
 b_i=\sum_{K_B}\sum_{K_A:\,|K_A\cap K_B|>0}
      \int_{K_A\cap K_B}N_i^B\phi_A^h\dd x
 \label{eq:galerkin-load-intersections}
\end{equation}
On each intersection, both expressions are fixed polynomials. These intersections form the common refinement and permit polynomial-exact quadrature: degree two for $P^1$/$P^1$ transfer and degree four for $P^2$/$P^2$ transfer on affine simplices~\cite{Dunavant1985,Keast1986}. Complete, nonoverlapping coverage is required. The intersection geometry and coupled solve distinguish projection from the point searches used by NN and BAR.

GP2 uses $P^1$ source and target spaces. GP3 constructs a $P^2$ source representation, including internally reconstructed edge values, and solves for a complete quadratic target field. Only original-vertex values are then supplied to the flow solver; the edge construction and the consequence of this restriction are detailed in Appendix~\ref{app:quadrature-mismatch}. On mixed-element meshes the spaces are $P^p$ on conforming simplicial subdivisions, not $Q^p$ on the original elements. Nominal approximation orders are two and three under smoothness and mesh-stability assumptions; accuracy after vertex restriction is assessed separately.

\subsection{Integral preservation, bounds, and the returned state}
\label{sec:global-correction}
\label{sec:galerkin-vertex-output}
\label{sec:galerkin-limiter}

When the unconstrained target space contains constants, Eq.~\eqref{eq:galerkin-orthogonality} preserves the integral of the complete finite-element field. It also preserves consistent coordinate moments represented by the test space. These identities do not imply conservation over arbitrary finite-volume cells, whose indicator functions generally lie outside a continuous finite-element space. Nor is the projection inverse a local operator: unchanged geometry in a subregion does not guarantee unchanged projected values there.

For GP2, consistent and solver-lumped integrals coincide only when $w_i^M=\int_\Omega N_i^M\dd x$ on both meshes. For GP3, source reconstruction and target vertex restriction introduce additional representation differences. An exact quadratic projection can acquire a second-order lumped-integral error after vertex restriction (Appendix~\ref{app:quadrature-mismatch}). A third-order pointwise approximation therefore does not imply a third-order conservation defect in the solver representation.

A separate global correction, denoted GC, can enforce Eq.~\eqref{eq:fv-discrete-integral} for any base transfer. For a raw target field $\phi_B^r$, the correction $c$ satisfies
\begin{equation}
 \sum_jw_j^Bc_j=Q_A(\phi_A)-Q_B(\phi_B^r),
 \qquad \ell_j\leq\phi_{B,j}^r+c_j\leq u_j
 \label{eq:global-correction-constraint}
\end{equation}
The defect is redistributed among nodes with available capacity, with saturated nodes fixed~\cite{ShashkovWendroff2004,KucharikShashkovWendroff2003}. The requested integral must lie within the weighted bounds. A small integral defect does not bound each correction, and bounds imposed separately on conserved components do not ensure positive pressure.

GP limiting uses source-vertex extrema in overlapping elements. The interval is widened by $\tau$ times the global source range, followed by clipping and redistribution preserving a selected lumped total. Hard limiting has $\tau=0$; soft limiting permits a finite excursion. Once active, limiting need not preserve Galerkin orthogonality, moments, or best approximation. Boundary enforcement and positivity corrections can change the field further. Conservation and accuracy are consequently assessed at specified stages, rather than assigned to the final state from properties of the unconstrained projection.

\section{Multistep response to a mesh change}
\label{sec:restart-consistency}

\subsection{Native and transferred histories}
\label{sec:unsteady-discretization}
\label{sec:native-transfer-decomposition}

Between mesh replacements the semi-discrete equations are
\begin{equation}
 M_M\dot U_M+F_M(U_M,t)=0
 \label{eq:semidiscrete}
\end{equation}
The state includes every transported variable required by the time integrator, including turbulence variables for RANS. Boundary constraints must be imposed consistently on the histories and on the equations advanced in time. At a constant time step, the target BDF2 equation is
\begin{equation}
 \Phi_B(Y;V,W)=\frac{M_B}{\Delta t}
        \left(\tfrac32Y-2V+\tfrac12W\right)+F_B(Y,t^{n+1})=0
 \label{eq:bdf2}
\end{equation}
Here $V$ and $W$ are the latest and older histories. The mesh is stationary during each time step. A continuously moving mesh would require compatible geometric histories and a discrete geometric conservation law~\cite{KoobusFarhat1999,GeuzaineGrandmontFarhat2003}; Eq.~\eqref{eq:bdf2} is not intended as that general formulation.

Let $U_A^k$ and $U_B^k$ denote native trajectories for the same physical problem on two meshes. The transferred histories and their target-referenced discrepancies are
\begin{equation}
 \widehat U_B^k=\calI_{AB}(U_A^k),\qquad
 \eps^k=\widehat U_B^k-U_B^k,\qquad k=n,n-1
 \label{eq:history-defects}
\end{equation}
The effective map $\calI_{AB}$ includes the processing required to form admissible histories. The discrepancy $\eps^k$ is not a pure interpolation error: it also contains the difference between the native source and target discretizations. Comparing with $U_B^k$ provides a common receiving evolution operator, not an exact continuum reference. The nonlinear decomposition in Appendix~\ref{app:reference-decomposition} makes this distinction explicit.

\subsection{Exact history forcing and conditional state response}

At a common trial state $Y$, changing only the two histories gives the exact identity
\begin{equation}
 \begin{aligned}
 &\Phi_B(Y;U_B^n+\eps^n,U_B^{n-1}+\eps^{n-1})
       -\Phi_B(Y;U_B^n,U_B^{n-1})\\
 &\hspace{2cm}=-\frac{M_B}{\Delta t}\eta^n,
 \qquad \eta^n=2\eps^n-\tfrac12\eps^{n-1}
 \end{aligned}
 \label{eq:history-residual-identity}
\end{equation}
No linearity of the spatial residual or small-error assumption is required. Matching the trial state, boundary treatment, and residual sign convention is essential.

The solved-state difference $\delta^{n+1}=\widehat U_B^{n+1}-U_B^{n+1}$ obeys a different relation. Let $r$ and $\widehat r$ be the algebraic residuals remaining after the native and perturbed physical steps. Subtraction and expansion about the native target state give
\begin{equation}
 \begin{aligned}
 A_B\delta^{n+1}&=M_B\eta^n+\Delta t(\widehat r-r)
                 -\Delta t\,\mathcal N_B(\delta^{n+1})\\
 A_B&=\tfrac32M_B+\Delta t J_B,
 \qquad J_B=\frac{\partial F_B}{\partial U}
 \end{aligned}
 \label{eq:linearized-bdf2-response}
\end{equation}
A locally Lipschitz Jacobian gives $\mathcal N_B(\delta)=O(\|\delta\|^2)$. If the step is locally invertible, the perturbation is small, and the algebraic errors are negligible, then
\begin{equation}
 \delta^{n+1}=R_B\eta^n+O(\Delta t\|\eta^n\|^2),
 \qquad R_B=A_B^{-1}M_B
 \label{eq:state-response}
\end{equation}
The bound is local to the mesh, state, time step, and norm. Limiter switching and incomplete inner convergence can invalidate this smooth-response approximation without invalidating Eq.~\eqref{eq:history-residual-identity}.

The steady-airfoil experiment places the same spatial discrepancy $\eps^\star=\calI_{AB}U_A^\star-U_B^\star$ in the two histories independently:
\begin{table}[htbp]
\centering\small
\caption{History perturbations used to separate the BDF2 contributions. The spatial discrepancy $\eps^\star$ is the same within each method; the native case supplies the receiving-mesh reference.}
\label{eq:naca-history-decomposition}
\begin{tabular}{lccc}
\toprule
Configuration & $\eps^n$ & $\eps^{n-1}$ & $\eta^n$\\
\midrule
Native & $0$ & $0$ & $0$\\
Both histories & $\eps^\star$ & $\eps^\star$ & $3\eps^\star/2$\\
Latest only & $\eps^\star$ & $0$ & $2\eps^\star$\\
Older only & $0$ & $\eps^\star$ & $-\eps^\star/2$\\
\bottomrule
\end{tabular}
\end{table}
Writing the corresponding responses as $\delta_f$, $\delta_n$, and $\delta_{n-1}$, local linearity predicts
\begin{equation}
 \delta_n\simeq\tfrac43\delta_f,\qquad
 \delta_{n-1}\simeq-\tfrac13\delta_f,\qquad
 \delta_f\simeq\delta_n+\delta_{n-1}
 \label{eq:naca-history-ratios}
\end{equation}
The older-only case tests whether a latest-state diagnostic misses a finite forcing. These collinear perturbations test directional response; they do not determine $R_B$ for arbitrary history fields.

\subsection{Aerodynamic response}
\label{sec:delayed-response}

For a differentiable target-mesh output $G_B$ with gradient $g_B^{n+1}$, Eq.~\eqref{eq:state-response} yields
\begin{equation}
 \Delta G^{n+1}=(g_B^{n+1})^TR_B\eta^n+O(\|\eta^n\|^2)
 \label{eq:force-response}
\end{equation}
The load response depends on the spatial discrepancy weighted by the target dynamics and output sensitivity, not just its integral. In particular, preserving the source integral imposes, componentwise,
\begin{equation}
 \bm1^TM_B\eps^k=Q_A(U_A^k)-Q_B(U_B^k)
 \label{eq:gc-native-mismatch}
\end{equation}
The right-hand side need not vanish because the native discrete states can have different integrals. Global correction therefore need not reduce the target-referenced state or load discrepancy.

After the first step, both the new solution difference and the remaining older history influence subsequent steps. A first-step load difference need not bound the later excursion. Disturbances generated away from the airfoil may produce a delayed response at an acoustic or convective scale, but neither the projection inverse nor an implicit step has compact spatial support. The numerical experiments therefore use finite-amplitude arrival thresholds rather than the first nonzero signal.


\section{Numerical experiments}
\label{sec:methodology}

\subsection{Static fields and repeated Euler transfer}
\label{sec:static-verification}
\label{sec:riemann-methodology}

Static transfer tests exclude flow-solver error. The smooth asymmetric field is
\begin{equation}
 \phi=\rho_\infty\left[2+0.3e^{0.8x/L_x+0.5y/L_y}
              \sin(2\pi x/L_x)\sin(2\pi y/L_y)\right]
 \label{eq:asymmetric-manufactured-field}
\end{equation}
Structured quadrilateral and independently generated mixed triangular/quadrilateral target meshes have nominal resolutions 64, 128, and 256; source resolution is half the target resolution. Relative nodal density errors are measured against the analytic field at target vertices. The effective spacing $h_{\mathrm{eff}}=N_{\mathrm{nodes}}^{-1/2}$ defines refinement ratios on the fixed domain. Pairwise orders are
\begin{equation}
 p_{k,k+1}=\frac{\log(E_k/E_{k+1})}{\log(h_k/h_{k+1})}
 \label{eq:pairwise-observed-order}
\end{equation}
A least-squares fit of $\log E$ against $\log h$ summarizes each three-level sequence. The discontinuous field has density 1.2 outside and 3.6 inside a circular inclusion of radius $0.25\min(L_x,L_y)$. The interface is not mesh-aligned. Unlimited and limited transfers are compared through overshoot, undershoot, and the source-to-target lumped-integral defect.

Repeated transfer is examined in a two-dimensional Euler Riemann calculation on $[0,5]\times[0,1]$. The baseline quadrilateral mesh has $501\times101$ vertices and spacing $h=0.01$. The gas is ideal with $\gamma=1.4$ and unit gas constant. The initial state is a shock tube at rest, discontinuous at $x=2$: $p=1$, $T=1$, $\rho=1$ for $x<2$ and $p=0.1$, $T=0.8$, $\rho=0.125$ for $x\geq2$, with zero velocity everywhere; these are the standard Sod states. The upper and lower boundaries are inviscid slip walls. The inlet and outlet planes are characteristic boundaries imposing total conditions and flow direction, $(p_0,T_0)=(1,1)$ at $x=0$ and $(0.1,0.8)$ at $x=5$, both with the direction $(1,0)$, so that each end holds the total state of the initial region adjacent to it. The calculation is a transient from this discontinuous state and is compared with a fixed-mesh calculation of the same problem, not with an exact Riemann solution. The calculation advances 400 physical steps with $\Delta t=10^{-3}$. Interior vertices receive fresh random displacements in $x$ relative to the original mesh; boundaries remain fixed, and geometric distortion does not accumulate between successive meshes. The nodal displacement bound is the smallest incident corner-triangle altitude divided by a safety factor, SF. SF2 permits twice the displacement of SF4.

Mesh replacement intervals of ten and two physical steps compare low- and high-frequency transfer. Both BDF2 histories are transferred at each mesh change. The low-frequency comparison contains 80 transfer evaluations; the high-frequency integral record contains 398. A fixed-mesh calculation supplies the reference. The comparison includes both transfer error and the change in spatial discretization; it is not an error against an exact Euler solution.

Density profiles are averaged in $y$ with a common sampling procedure. The transition position is its 50\% crossing, and its width is the distance between the 10\% and 90\% crossings. Changes from the fixed-mesh reference are expressed in baseline cells. These metrics describe a resolved density transition rather than presume an independently identified sharp shock. Density and energy integral changes are evaluated immediately across each transfer, separately from physical evolution. NN1, NN1+GC, BAR2, and GP2 are compared across three random deformation realizations at both amplitudes; other methods provide a broader SF4 frequency comparison.

\subsection{Euler airfoil and controlled BDF2 histories}
\label{sec:naca-fixed-methodology}

Euler flow over NACA~0012 is calculated at $M_\infty=0.1$ and $\alpha=6^\circ$, with $p_\infty=\rho_\infty=T_\infty=1$, $a_\infty=\sqrt{1.4}$, $U_\infty=0.118322$, and $\Delta t=10^{-3}$. The source and independent target meshes contain 16,255 and 22,508 vertices, respectively, with 32,077 and 44,583 triangles. Both retain the same 370 airfoil vertices and 63-vertex farfield polygon at $100c$, so the discretized surface is unchanged. Forces are pressure-only, with unit reference chord and area and a moment reference at $x/c=0.25$. Nonzero steady drag is treated as a numerical load, not skin friction or shock-induced drag.

Each native steady solution is continued with equal initial BDF2 histories. The history combinations in Table~\ref{eq:naca-history-decomposition} are then formed on the target mesh and compared with its native continuation at the same physical step. State differences use the nondimensional volume-weighted norm $\|v\|_{M_B}=(v^TM_Bv)^{1/2}$ over the stored mean-flow components. Supplementary tests also report the unweighted Euclidean norm $\|v\|_2$; these norms weight the near-body and large farfield cells differently. Relative superposition defects are
\begin{equation}
 e_{\mathrm{sup}}^U=
 \frac{\|\delta_f-\delta_n-\delta_{n-1}\|_{M_B}}{\|\delta_f\|_{M_B}},
 \qquad
 e_{\mathrm{sup}}^{C_L}=
 \frac{|\Delta C_{L,f}-\Delta C_{L,n}-\Delta C_{L,n-1}|}{|\Delta C_{L,f}|}
 \label{eq:history-superposition-error}
\end{equation}
State-norm ratios are nonnegative; force ratios retain their sign.

The solver reports $-\Phi_B$. Its initial residuals are compared before an update using two common-trial pairs: older-only minus native at $U_B^\star$, and latest-only minus both-histories at $\widehat U_B^\star$. Their differences in $\eta$ are $-\eps^\star/2$ and $+\eps^\star/2$, respectively. Double-precision residuals and a wall-normal/tangential split separate temporal forcing from slip-wall enforcement. Additional cancellation and amplitude tests are described in Appendix~\ref{app:history-checks}.

A repeated-remeshing experiment perturbs vertices at wall distances between $0.01c$ and $20c$, while retaining the boundaries. NN2 and GP2 use the same two deformation realizations, with 39 mesh replacements followed by evolution in a 400-step interval. Load differences are formed against the corresponding fixed-mesh trajectory. Slopes per 100 physical steps are descriptive finite-window measures, not estimates of a long-time drift or independent-sample uncertainty.

A separate single-transfer experiment perturbs narrow bands beginning at wall distances $0.05c$, $0.10c$, and $0.20c$. Fixed pressure shells measure the root-mean-square pressure difference, and force onset is defined by a fraction of the response maximum. For a shell at $r$ inward of a band at $d$, the low-Mach acoustic scale is $(d-r)/(a_\infty\Delta t)\simeq845(d-r)$ steps. Immediate disturbances extend up to approximately $0.08c$ inward of the displaced band; those shells are excluded from the propagation fit. Results are interpreted as threshold-dependent propagation, with a common fixed-mesh reference, rather than strict locality of interpolation.

\subsection{Fixed-body RANS and numerical controls}
\label{sec:rans-methodology}

The RANS case retains $M_\infty=0.1$ and $\alpha=6^\circ$, with $Re_c=6\times10^6$ and the Spalart--Allmaras (SA) and shear-stress-transport (SST) closures~\cite{SpalartAllmaras1994,Menter1994}. The production mesh has 23,271 vertices, 9,686 quadrilaterals in the wall layer and wake region, 26,851 triangles, and 256 airfoil vertices. The farfield is at $100c$. The first off-wall spacing is approximately $5\times10^{-6}c$, and the quadrilateral region extends to about $0.11c$. This is a remeshing experiment about a nominally steady viscous baseline, not a turbulence-model validation.

The SU2 calculations use HLLC mean-flow fluxes with MUSCL reconstruction and scalar upwinding without MUSCL for turbulence transport. BDF2 physical steps are solved by dual-time iteration with CFL~6 and FGMRES/ILU linear solves. The production cap is 800 inner iterations, with a density-residual target $\log_{10}(\mathrm{rms}\,\rho)=-12$ checked after ten iterations. The dimensional physical step $\Delta t_0=10^{-3}$~s corresponds to $\Delta t_0^*=0.347224$ under the chosen normalization. With $U_\infty^*=0.1$, the 400-step interval spans approximately 13.9 chord flow-through times. A first-order startup step is ten times smaller. Each closure has a no-remeshing reference continued for 800 steps from matching histories.

The steady SST solution reaches the density-residual target. The production SA steady calculation reaches its 20,000-iteration cap at $-11.2$, with lift and drag stationary to $10^{-7}$ over the final 2,000 iterations. The intermediate-mesh SST baseline stops at approximately $-11.1$. These residual levels and the reference continuations are distinguished from convergence of the remeshed physical steps.

The same prescribed mesh sequences are used for NN2 and limited GP2. At each replacement, 12,602 vertices may move and 10,669 are fixed. No vertex with solver wall distance below $0.01c$ moves; the geometric mask uses distance to the nearest airfoil vertex, and part of the outer quadrilateral layer can move. SF4 limits displacements, and a minimum corner-area ratio of 0.2 preserves valid cells. The boundaries remain fixed. Mesh replacement occurs every ten physical steps, independently of the evolving solution.

Mean-flow and turbulence histories are transferred at both required time levels. SA transports $\widetilde\nu$; SST transfers density-weighted $k$ and $\omega$, followed by primitive-variable positivity floors. Wall distance and eddy viscosity are recomputed. Turbulence integrals are not invariants of their subsequent transport equations. The SST transfer-integral diagnostic is evaluated before the final positivity floor and does not, by itself, establish conservation of the fully processed restart.

For output $Q$, define the case-minus-reference response and paired method difference by
\begin{equation}
 D_{m,Q}(t)=C_{Q,m}(t)-C_{Q,\mathrm{ref}}(t),
 \qquad P_Q(t)=D_{\mathrm{NN2},Q}(t)-D_{\mathrm{GP2},Q}(t)
 \label{eq:rans-comparisons}
\end{equation}
Pressure and friction contributions use identical surface quadrature for case and reference. Over common physical sample times, write $\mathcal R(q)=(N^{-1}\sum_i q(t_i)^2)^{1/2}$. Sensitivity to a control setting $c$ is measured by
\begin{equation}
 S(q;c)=\frac{\mathcal R(q_c-q_0)}{\mathcal R(q_0)}
 \label{eq:control-sensitivity}
\end{equation}
This measures the difference between histories, not the fractional change of their RMS amplitudes. The latter is $[\mathcal R(q_c)-\mathcal R(q_0)]/\mathcal R(q_0)$.

Controls cover a common 38-sample window and three mesh events from one realization. Inner-iteration caps are increased to 1600 for both closures and to 6400 for SST. Time-step controls use $\Delta t_0/2$ for both closures and $\Delta t_0/4$ for SA, with identical displacement fields applied at matching physical times and a separate no-remeshing reference for each step size. Two finer steady meshes assess sensitivity of the absolute loads. A short SA remeshing comparison on the intermediate mesh uses the same displacement rules but a different node-wise realization; it is treated as cross-mesh robustness rather than isolated spatial convergence.


\section{Results}
\label{sec:results}

\subsection{Interpolation order verification}
\label{sec:results-asymmetric}

Figure~\ref{fig:asym-l2-convergence} shows the smooth-field refinement results. NN1 is first order; NN2, BAR2, and GP2 are consistent with second order; NN3 is consistent with third order. BAR3 and GP3 have fitted slopes above three over the three-level interval, approaching approximately 3.2 on the finest pair. This is finite-range convergence behavior, not evidence of a higher designed order. On the finest mixed grid, the relative errors are approximately $2.6\times10^{-7}$ for BAR3, $3.4\times10^{-7}$ for GP3, and $2.6\times10^{-6}$ for NN3. These differences resolve smooth interpolation performance on the tested mesh family; they do not predict a corresponding ranking of restarted flow solutions.

\begin{figure}[htbp]
    \centering
    \begin{tikzpicture}
        \begin{groupplot}[
            group style={group size=2 by 1,horizontal sep=1.2cm},
            transfer convergence axis,
            width=0.45\textwidth,
            height=0.34\textwidth,
        ]
        \nextgroupplot[
            title={Structured quadrilateral grid},
            ylabel={Relative $L_2$ error},
            legend to name=legend:asym-l2,
            legend columns=4,
        ]
        \addplot[scheme NN1] table[col sep=comma,x=h_Struct,y=Struct_NN1]{\GeneratedDir/asymmetric_convergence_L2.csv}; \addlegendentry{NN1}
        \addplot[scheme NN2] table[col sep=comma,x=h_Struct,y=Struct_NN2]{\GeneratedDir/asymmetric_convergence_L2.csv}; \addlegendentry{NN2}
        \addplot[scheme NN3] table[col sep=comma,x=h_Struct,y=Struct_NN3]{\GeneratedDir/asymmetric_convergence_L2.csv}; \addlegendentry{NN3}
        \addplot[scheme BAR2] table[col sep=comma,x=h_Struct,y=Struct_BAR2]{\GeneratedDir/asymmetric_convergence_L2.csv}; \addlegendentry{BAR2}
        \addplot[scheme BAR3] table[col sep=comma,x=h_Struct,y=Struct_BAR3]{\GeneratedDir/asymmetric_convergence_L2.csv}; \addlegendentry{BAR3}
        \addplot[scheme GP2] table[col sep=comma,x=h_Struct,y=Struct_GP2]{\GeneratedDir/asymmetric_convergence_L2.csv}; \addlegendentry{GP2}
        \addplot[scheme GP3] table[col sep=comma,x=h_Struct,y=Struct_GP3]{\GeneratedDir/asymmetric_convergence_L2.csv}; \addlegendentry{GP3}
        \addplot[reference slope] coordinates {(0.014,8.0e-3) (0.007,4.0e-3)} node[pos=0.55,above,sloped]{\scriptsize $h$};
        \addplot[reference slope] coordinates {(0.014,1.2e-3) (0.007,3.0e-4)} node[pos=0.55,above,sloped]{\scriptsize $h^2$};
        \addplot[reference slope] coordinates {(0.014,1.2e-4) (0.007,1.5e-5)} node[pos=0.55,above,sloped]{\scriptsize $h^3$};

        \nextgroupplot[title={Unstructured mixed grid}]
        \addplot[scheme NN1] table[col sep=comma,x=h_Mixed,y=Mixed_NN1]{\GeneratedDir/asymmetric_convergence_L2.csv};
        \addplot[scheme NN2] table[col sep=comma,x=h_Mixed,y=Mixed_NN2]{\GeneratedDir/asymmetric_convergence_L2.csv};
        \addplot[scheme NN3] table[col sep=comma,x=h_Mixed,y=Mixed_NN3]{\GeneratedDir/asymmetric_convergence_L2.csv};
        \addplot[scheme BAR2] table[col sep=comma,x=h_Mixed,y=Mixed_BAR2]{\GeneratedDir/asymmetric_convergence_L2.csv};
        \addplot[scheme BAR3] table[col sep=comma,x=h_Mixed,y=Mixed_BAR3]{\GeneratedDir/asymmetric_convergence_L2.csv};
        \addplot[scheme GP2] table[col sep=comma,x=h_Mixed,y=Mixed_GP2]{\GeneratedDir/asymmetric_convergence_L2.csv};
        \addplot[scheme GP3] table[col sep=comma,x=h_Mixed,y=Mixed_GP3]{\GeneratedDir/asymmetric_convergence_L2.csv};
        \addplot[reference slope] coordinates {(0.012,8.0e-3) (0.006,4.0e-3)} node[pos=0.55,above,sloped]{\scriptsize $h$};
        \addplot[reference slope] coordinates {(0.012,1.2e-3) (0.006,3.0e-4)} node[pos=0.55,above,sloped]{\scriptsize $h^2$};
        \addplot[reference slope] coordinates {(0.012,1.2e-4) (0.006,1.5e-5)} node[pos=0.55,above,sloped]{\scriptsize $h^3$};
        \end{groupplot}
    \end{tikzpicture}
    \vspace{0.5em}
    \ref{legend:asym-l2}
    \caption{Convergence of the density-transfer error for the smooth asymmetric field. Curves show the relative nodal $L_2$ error as the mesh is refined; gray guides correspond to $O(h)$, $O(h^2)$, and $O(h^3)$.}
    \label{fig:asym-l2-convergence}
\end{figure}

\subsection{Boundedness and conservation}
\label{sec:results-boundedness}

The circular-inclusion test separates boundedness from integral preservation. Table~\ref{tab:disc-selected} compares representative transfers at the middle structured resolution, $64\to128$. Unlimited NN2 exceeds the density plateaus by 0.8031 above and 0.4157 below. BAR3 reduces these excursions but remains unbounded. Barth--Jespersen limiting removes the measured new extrema; where it suppresses the NN2 reconstruction, however, the integral defect increases toward the NN1 value. Global correction then restores the integral without reintroducing an extremum. Hard-limited GP2 is bounded and conservative in this representation, whereas vertex-only GP3 retains a nonzero lumped-integral defect. The methods therefore cannot be characterized by one combined measure of order, boundedness, and conservation.

\begin{table}[htbp]
\centering\small
\caption{Representative discontinuous-field transfers. Overshoot and undershoot are excursions beyond the exact density range $[1.2,3.6]$; the last column is the source-to-target defect in the solver's lumped density integral. BJ denotes Barth--Jespersen limiting and GC the bounded global conservation correction.}
\label{tab:disc-selected}
\begin{tabular}{lrrr}
\toprule
Transfer & Overshoot & Undershoot & Integral defect \\
\midrule
NN2, unlimited & 0.8031 & 0.4157 & $-5.26\times10^{-5}$ \\
NN2, BJ & 0 & 0 & $-4.46\times10^{-3}$ \\
NN2, BJ+GC & 0 & 0 & $2.95\times10^{-14}$ \\
BAR3, unlimited & 0.1766 & 0.1811 & $-8.80\times10^{-7}$ \\
BAR3, BJ+GC & 0 & 0 & $5.59\times10^{-14}$ \\
GP2, hard limited & 0 & 0 & $3.44\times10^{-15}$ \\
GP3, hard limited, vertex output & 0 & 0 & $-4.83\times10^{-7}$ \\
\bottomrule
\end{tabular}
\end{table}

For a separate smooth transfer, global correction reduces the density-integral defect to order $10^{-15}$ while changing the density $L_1$ error by at most 0.8\%. Its one-step accuracy penalty is small in that test. Whether the same scalar constraint controls repeated spatial distortion is examined next.

\subsection{Accumulated broadening under repeated Euler remeshing}
\label{sec:results-riemann}

At the low-frequency setting, every method locates the density transition within 0.14 baseline cell of the fixed-mesh reference. Width changes are also small: 0.57 cell for NN1, 0.25 for BAR2, and at most 0.01 in magnitude for the remaining variants. More frequent transfer changes the width much more than the position (Table~\ref{tab:riemann-stress-new}). At SF4, NN1 broadening grows to 21.41 cells and BAR2 broadening to 3.24 cells, while GP2 has no width change to the reported precision. At SF2, NN1 and BAR2 broaden by 26.36 and 22.91 cells, respectively, compared with 0.06 for GP2.

\begin{table}[htbp]
\centering\small
\caption{Broadening of the density transition relative to the fixed-mesh reference. Width increases and the largest absolute position error among the listed methods are expressed in baseline cells ($h=0.01$). GC is the global correction; SF2 doubles the displacement bound of SF4.}
\label{tab:riemann-stress-new}
\begin{tabular}{lrrrrr}
\toprule
Condition & NN1 width & NN1+GC width & BAR2 width & GP2 width & max. $|\Delta x|/h$ \\
\midrule
80 transfers, SF4  & 0.57 & --    & 0.25 & 0.00 & 0.14 \\
398 transfers, SF4 & 21.41 & 18.36 & 3.24 & 0.00 & 0.15 \\
398 transfers, SF2 & 26.36 & 25.65 & 22.91 & 0.06 & 0.91 \\
\bottomrule
\end{tabular}
\end{table}

The hierarchy repeats across three random realizations. At SF4, excess widths range from 21.41 to 22.89 cells for NN1, 3.12 to 3.24 for BAR2, and 0 to 0.01 for GP2. At SF2 the respective ranges are 24.98--27.14, 22.90--22.98, and 0.03--0.06 cells. Position errors remain below 0.5 cell for these uncorrected comparisons. The dominant degradation is thus broadening rather than systematic translation of the feature.

Global correction isolates the distinction between integral and spatial error. Across all corrected NN1 realizations, the absolute cumulative density-integral defect does not exceed $9.4\times10^{-14}$, yet excess width remains 18.36--21.28 cells at SF4 and 23.80--26.32 cells at SF2. Correcting the total does not recover the lost profile. The result does not establish a unique benefit of Galerkin projection: NN2, NN3, BAR3, and GP3 also have much smaller width changes than NN1 in the SF4 comparison. Appendix~\ref{app:riemann-support} retains those alternatives and the complete realization data.

\subsection{Response to the two BDF2 histories}
\label{sec:results-bdf2-new}

For the independent Euler airfoil meshes, second- and third-order transfers give target-referenced state discrepancies near $7.0\times10^{-4}$, with less than 3\% spread; NN1 gives $1.47\times10^{-3}$. This clustering is not evidence that the higher-order interpolants are equally accurate. The discrepancy includes the difference between native discrete equilibria, whereas the static verification compares against a known field.

The history experiment yields a clearer temporal result. Figure~\ref{fig:naca-history-new} shows a finite older-only lift response for every transfer, although the latest history is exact. For GP2, the older-only lift difference is $+2.539\times10^{-3}$, compared with $-7.575\times10^{-3}$ when both histories are perturbed. The latest-state discrepancy alone therefore misses a nonzero forcing of the next physical step.

\begin{figure}[htbp]
\centering
\begin{tikzpicture}
\begin{axis}[
 ybar,
 width=0.82\textwidth,height=5.8cm,
 ylabel={$\Delta C_L$ after the first BDF2 step},
 symbolic x coords={NN1,NN2,BAR2,GP2,GP3},xtick=data,
 bar width=5pt,
 legend style={font=\small,draw=none,at={(0.5,1.03)},anchor=south,legend columns=3,/tikz/every even column/.append style={column sep=8pt}},
 legend cell align=left,
 grid=major,
 scaled y ticks=false,
]
\addplot table[col sep=comma,x=method,y=full]{\DataDir/naca_first_step_CL_plot.csv};\addlegendentry{Both histories}
\addplot table[col sep=comma,x=method,y=latest]{\DataDir/naca_first_step_CL_plot.csv};\addlegendentry{Latest only}
\addplot table[col sep=comma,x=method,y=older]{\DataDir/naca_first_step_CL_plot.csv};\addlegendentry{Older only}
\end{axis}
\end{tikzpicture}
\caption{Lift change after the first physical step on the receiving Euler airfoil mesh, relative to its native continuation. Within each method the same spatial discrepancy is assigned to both histories, the latest history only, or the older history only. Opposite signs of the older-only and latest-only responses follow the BDF2 history weights.}
\label{fig:naca-history-new}
\end{figure}

Across all seven transfer variants, the latest/full state-norm ratio is 1.3330--1.3333 and the older/full magnitude is 0.3335--0.3341, close to the BDF2 weights in Eq.~\eqref{eq:naca-history-ratios}. Volume-weighted state superposition defects are 1.8--3.5\%, and lift superposition defects are 0.28--1.49\%. These measurements support the predicted directional response, not an independently validated linear solution operator for arbitrary perturbations.

The common-trial residual comparison identifies an additional boundary effect. NN and BAR satisfy the history-forcing identity to approximately $10^{-7}$ relative error. For GP2 and GP3, the larger whole-vector mismatch is confined to wall-normal momentum at the 370 slip-wall vertices. The normal mismatch equals the corresponding transferred forcing, while tangential differences are of order $10^{-11}$. At the projected boundary rows, the same boundary operation must be applied to the predicted history forcing; the interior and tangential components then recover the identity. The distinction is between the unconstrained projected history and the history admitted by the receiving equations, not a failure of the BDF2 identity.

Supplementary cancellation tests produce very small first-step force differences, but not comparably small volume-weighted state differences. Reducing the imposed amplitude decreases Euclidean-state and lift superposition remainders, whereas the relative volume-weighted remainder increases. Both diagnostics are retained in Appendix~\ref{app:history-checks}; neither an iterative floor nor a uniform quadratic remainder is inferred from these observations.

Global correction changes the NN1 full-history lift response from $5.7873\times10^{-3}$ to $5.7924\times10^{-3}$ and the BAR2 response from $-7.2852\times10^{-3}$ to $-7.2850\times10^{-3}$. In these cases the integral constraint makes little difference to the load response, consistent with the spatial weighting in Eq.~\eqref{eq:force-response}.

\subsection{Accumulation and delayed Euler load response}
\label{sec:results-euler-airfoil-new}

The repeated Euler airfoil calculation shows a method-dependent finite-window drag trend. The GP2 slopes are $6.84\times10^{-6}$ and $6.38\times10^{-6}$ per 100 physical steps for the two realizations; NN2 gives $3.95\times10^{-5}$ and $3.84\times10^{-5}$. The ratio of the two-realization means is approximately 5.9. This is a repeatable contrast for the tested discretization and interval, not a universal accuracy or cost ranking. The 400-step Euler interval is only about 0.047 chord flow-through times, and lift trends are less consistent between realizations.

The localized-transfer experiment also shows why the first-step load difference is insufficient. On the longer $0.20c$-band case, pressure-shell propagation speeds inferred from 10--75\% amplitude thresholds range from 1.09 to 1.27, bracketing $a_\infty=1.183$. This range measures detector sensitivity for one event, not an ensemble confidence interval. It supports an acoustic-scale delayed response outside the immediate transfer footprint. Separately, the NN1 lift difference peaks three steps after the first post-transfer sample, at 4.66 times its initial value; the sampled second- and third-order responses instead peak at the first step and relax thereafter.

\subsection{RANS transfer integrals and aerodynamic resolution}
\label{sec:results-rans-new}

With turbulence transport, GP2 again preserves the monitored integrals more closely than NN2 (Table~\ref{tab:rans-integrals-new}). Maximum density and total-energy defects differ by roughly two orders of magnitude. For the turbulence inventories the NN2/GP2 ratios are approximately 247 for SA and 37 for both SST quantities. These are instantaneous transfer defects at the measurement stage described in Sec.~\ref{sec:rans-methodology}, before the final SST positivity floor; they are not conservation errors accumulated through the turbulence equations.

\begin{table}[htbp]
\centering\small
\caption{Largest absolute source-relative integral change across the repeated RANS history transfers. Each column is maximized separately over the methods' two realizations. The SST values precede the final positivity floor. A larger NN2/GP2 ratio indicates better preservation of that integral by GP2 at this stage.}
\label{tab:rans-integrals-new}
\pgfplotstabletypeset[publication table,
 columns={closure,quantity,NN2_max,GP2_max,ratio},
 columns/closure/.style={string type,column name={Closure}},
 columns/quantity/.style={string type,column name={Inventory}},
 columns/NN2_max/.style={sci,precision=2,column name={NN2}},
 columns/GP2_max/.style={sci,precision=2,column name={GP2}},
 columns/ratio/.style={fixed,precision=0,column name={NN2/GP2}}
]{\GeneratedDir/reassessed_rans_integrals.csv}
\end{table}

The load comparison is less separated. Over the repeated-remeshing interval, total-drag RMS deviations from the no-remeshing references range from $6.07\times10^{-5}$ to $6.79\times10^{-5}$. Paired NN2--GP2 differences have RMS $0.77\times10^{-5}$ to $1.11\times10^{-5}$, and the ordering of the methods changes with closure and realization. Similar aggregate deviations do not imply identical histories or aerodynamic equivalence.

The temporal controls distinguish the overall disturbance from that smaller paired difference (Table~\ref{tab:rans-temporal-main}). Halving the step changes the time-aligned individual lift histories by 15.1--16.5\% of their production RMS for SA and 19.9--20.5\% for SST; drag changes are 8.0--10.2\% and 10.3--10.8\%, respectively. Their RMS amplitudes change much less. In contrast, the paired lift and drag differences change by 75.2\% and 128.7\% for SA, and by 146.4\% and 110.3\% for SST. The SA quarter-step comparison increases the change of the paired difference further. Thus the overall disturbance remains of comparable magnitude, but neither the detailed transient nor a method ranking is established as time-step converged.

\begin{table}[htbp]
\centering\small
\caption{Time-step sensitivity at common physical times over the three-event RANS control window. Values are $100S$ from Eq.~\eqref{eq:control-sensitivity}. Individual columns give the larger NN2 or GP2 sensitivity; paired columns refer to their difference. These are RMS differences between histories, not percentage changes in their RMS amplitudes.}
\label{tab:rans-temporal-main}
\begin{tabular}{llrrrr}
\toprule
 & & \multicolumn{2}{c}{Individual histories (\%)} & \multicolumn{2}{c}{Paired difference (\%)}\\
\cmidrule(lr){3-4}\cmidrule(lr){5-6}
Closure & Step & $C_L$ & $C_D$ & $C_L$ & $C_D$\\
\midrule
\ControlRow{SA}{halfstep}{$\Delta t_0/2$}
\ControlRow{SA}{quarterstep}{$\Delta t_0/4$}
\ControlRow{SST}{halfstep}{$\Delta t_0/2$}
\bottomrule
\end{tabular}
\end{table}

SA reaches the density-residual target at every reduced-step control point. Its production-step cap extension has a smaller, but nonzero, effect on the paired difference. SST is less conclusive: most moved-grid control steps still reach the cap when it is raised from 1600 to 6400, and the half-step run is also incompletely converged. Its observed time-step sensitivity therefore includes possible iterative-error contributions. Appendix~\ref{app:rans-controls} quantifies these limits, including the more sensitive friction contribution.

Three steady viscous mesh levels show sensitivity of the absolute loads, but do not establish an asymptotic refinement sequence. That comparison is not an error bound on the paired NN2--GP2 difference on a common mesh. Likewise, the finer-mesh remeshing window changes the realized nodal disturbance and cannot isolate mesh convergence of the transient. Within these limits, RANS extends the integral comparison to turbulence variables and viscous flow, while providing no resolved aerodynamic advantage for either transfer method.

\section{Conclusions}
\label{sec:conclusions}

Mesh-to-mesh transfer has been assessed through its approximation properties, the BDF2 histories it supplies, and the resulting aerodynamic response. The analysis distinguishes transfer consistency from the native source--target discretization difference, while numerical controls assess selected sensitivities of the receiving calculation.

Global conservation alone does not control repeated spatial distortion. Across three Euler deformation realizations, correcting the density integral to roundoff leaves large transition broadening. Higher-order pointwise interpolation and Galerkin projection both perform well in parts of the comparison, so the result cannot be attributed uniquely to finite-volume local conservation.

The controlled airfoil restarts establish that both BDF2 histories matter: an exact latest history does not prevent a next-step disturbance when the older history is perturbed. Directional response ratios follow the temporal coefficients, subject to finite state and force superposition errors. Boundary enforcement modifies the projected momentum as well, and must be included in the definition of the accepted history. Small first-step forces do not certify a small state discrepancy or bound the later load excursion.

The aerodynamic benefit of integral preservation is not universal in the tested calculations. GP2 reduces the finite-window Euler drag trend relative to NN2, while the RANS comparisons show much smaller integral defects without a time-step-resolved load ranking. The time-step controls preserve the overall scale of the disturbance more closely than the small difference between methods; SST remains additionally limited by incomplete inner convergence. These results support separate assessment of accuracy, admissibility, integral changes, history consistency, and output response. They do not establish aerodynamic equivalence, a universally preferred transfer, or convergence for moving-domain problems.


\appendix
\section{Representation and reference-state identities}
\label{app:quadrature-mismatch}

GP3 reconstructs the source value at the midpoint of edge $(i,j)$ as
\begin{equation}
 \phi_{ij}=\tfrac12(\phi_i+\phi_j)
       -\tfrac18d_{ij}^T\overline H_{ij}d_{ij},
 \qquad d_{ij}=x_i-x_j,\qquad \overline H_{ij}=\tfrac12(H_i+H_j)
 \label{eq:gp3-edge-value}
\end{equation}
For a quadratic field, the curvature term is exactly the difference between the endpoint average and midpoint value. Quadratic projection can consequently be exact even though restriction changes its solver integral. For example, $f(x)=x^2$ on a uniform partition of $[0,1]$ belongs to the quadratic finite-element space, but integration of only its vertex values with linear nodal weights gives
\begin{equation}
 Q_h(f)=\frac13+\frac{h^2}{6}
 \label{eq:vertex-restriction-counterexample}
\end{equation}
The second-order integral discrepancy arises solely from the representation change, not an error in the complete projection.

\subsection{Discrepancy relative to a native target state}
\label{app:reference-decomposition}

Let $\Pi_Mu^k$ represent a common exact field in the solver representation, with native error $e_M^k=U_M^k-\Pi_Mu^k$. For the complete, possibly nonlinear, transfer map,
\begin{align}
 \eps^k={}&\underbrace{\calI_{AB}(\Pi_Au^k)-\Pi_Bu^k}_{\text{transfer consistency}}\notag\\
 &+\underbrace{\calI_{AB}(U_A^k)-\calI_{AB}(\Pi_Au^k)-e_B^k}_{\text{source and target discretization contributions}}
 \label{eq:defect-decomposition}
\end{align}
Only a linear map permits the second line to be written $\calI_{AB}e_A^k-e_B^k$. A load difference likewise separates as
\begin{align}
 G_B(\widehat U_B^k)-G_A(U_A^k)
 ={}&G_B(U_B^k)-G_A(U_A^k)\notag\\
 &+G_B(\widehat U_B^k)-G_B(U_B^k)
 \label{eq:observable-decomposition}
\end{align}
The first difference is native mesh sensitivity; the second is the target-referenced restart response. Neither alone is an independently known continuum error.

\section{Repeated Euler transfer comparisons}
\label{app:riemann-support}

Table~\ref{tab:riemann-frequency-app} retains the higher-order alternatives to the representative main-text comparison. Table~\ref{tab:riemann-seeds-app} gives all three realizations, including NN1 with global correction. Transfer-integral defects and profile changes are separate measures: cancellation between signed integral changes does not bound local distortion.

\begin{table}[htbp]
\centering\small
\caption{SF4 frequency comparison for all transfer variants. $\Delta n$ is the number of physical steps between mesh changes; $\Delta x/h$ and $\Delta w/h$ are signed position and width changes relative to the fixed-mesh reference.}
\label{tab:riemann-frequency-app}
\pgfplotstabletypeset[publication table,
 columns={scheme,dIters,position_error_cells,width_change_cells,cumulative_density_change},
 columns/scheme/.style={string type,column name={Transfer}},
 columns/dIters/.style={fixed,precision=0,column name={$\Delta n$}},
 columns/position_error_cells/.style={fixed,precision=2,column name={$\Delta x/h$}},
 columns/width_change_cells/.style={fixed,precision=2,column name={$\Delta w/h$}},
 columns/cumulative_density_change/.style={sci,precision=2,column name={$\sum\Delta I_\rho$}}
]{\DataDir/riemann_frequency_full.csv}
\end{table}

\begin{table}[htbp]
\centering\small
\caption{High-frequency Euler comparison across three random deformation realizations. GC identifies the corrected NN1 calculations; SF2 doubles the SF4 displacement bound. Position and width changes are in baseline cells. The final column is the cumulative density-integral change introduced by transfer.}
\label{tab:riemann-seeds-app}
\pgfplotstabletypeset[publication table,
 columns={scheme,global_correction,safe_factor,seed,position_error_cells,width_change_cells,cumulative_density_change},
 columns/scheme/.style={string type,column name={Transfer}},
 columns/global_correction/.style={string type,column name={GC},string replace*={Global}{on}},
 columns/safe_factor/.style={fixed,precision=0,column name={SF}},
 columns/seed/.style={fixed,precision=0,column name={Realization}},
 columns/position_error_cells/.style={fixed,precision=2,column name={$\Delta x/h$}},
 columns/width_change_cells/.style={fixed,precision=2,column name={$\Delta w/h$}},
 columns/cumulative_density_change/.style={sci,precision=1,column name={$\sum\Delta I_\rho$}}
]{\DataDir/riemann_amplitude_seeds.csv}
\end{table}
\FloatBarrier

\section{Controlled-history diagnostics}
\label{app:history-checks}

\subsection{Directional response and boundary enforcement}

Table~\ref{tab:naca-superposition-app} gives the complete response ratios. The state ratios use the volume-weighted norm; the older-only force ratio is signed. Their agreement with the temporal coefficients is approximate and does not remove the superposition errors.

\begin{table}[htbp]
\centering\footnotesize
\caption{Euler BDF2 response ratios and relative superposition defects. Subscripts $f$, $n$, and $n-1$ denote both-histories, latest-only, and older-only perturbations. Ideal directional ratios are $4/3$ and $-1/3$; the state columns show their nonnegative norm ratios.}
\label{tab:naca-superposition-app}
\pgfplotstabletypeset[publication table,
 columns={method,state_superposition_defect,state_latest_full,state_older_full,CL_superposition_defect,CL_latest_full,CL_older_full},
 columns/method/.style={string type,column name={Transfer}},
 columns/state_superposition_defect/.style={sci,precision=1,column name={$e_{\mathrm{sup}}^U$}},
 columns/state_latest_full/.style={fixed,precision=4,column name={$\|\delta_n\|_M/\|\delta_f\|_M$}},
 columns/state_older_full/.style={fixed,precision=4,column name={$\|\delta_{n-1}\|_M/\|\delta_f\|_M$}},
 columns/CL_superposition_defect/.style={sci,precision=1,column name={$e_{\mathrm{sup}}^{C_L}$}},
 columns/CL_latest_full/.style={fixed,precision=4,column name={$\Delta C_{L,n}/\Delta C_{L,f}$}},
 columns/CL_older_full/.style={fixed,precision=4,column name={$\Delta C_{L,n-1}/\Delta C_{L,f}$}}
]{\DataDir/naca_superposition.csv}
\end{table}

In the common-trial residual test, the normal momentum mismatch is $2.125\times10^{-6}$ for GP2 and $6.066\times10^{-7}$ for GP3, equal to the normal history forcing at the slip wall. The tangential mismatch is of order $10^{-11}$. NN and BAR preserve wall-normal momentum much more closely in this particular transfer. The comparison must apply the same boundary operation to both the measured residual and the predicted history forcing; the raw unconstrained projection is a different state.

\subsection{Cancellation of the initial history forcing}

For $\eps^n=e/4$ and $\eps^{n-1}=e$, the initial combination $\eta^n$ is zero. At a common trial state the two BDF2 equations are then identical, and an exact solve on the same local solution branch should give the same next state. Table~\ref{tab:naca-cancellation-app} separates the measured force and state diagnostics. First-step lift differences are of order $10^{-9}$, but relative state differences are approximately $2\times10^{-4}$ in the Euclidean norm and 0.043--0.090 in the volume-weighted norm. The small force response therefore does not establish cancellation of the complete numerical state, and no independently measured iterative floor is assigned to either norm.

\begin{table}[htbp]
\centering\small
\caption{Initial history-cancellation test. Both histories are non-native but $2\eps^n-\eps^{n-1}/2=0$. First-step force differences and state differences use the same native continuation; both state norms are shown to avoid interpreting a small load difference as complete state agreement.}
\label{tab:naca-cancellation-app}
\pgfplotstabletypeset[publication table,
 columns={method,dCL_step2,dCD_step2,state_diff_step2_over_eps,state_diff_step2_over_eps_volume_weighted},
 columns/method/.style={string type,column name={Transfer}},
 columns/dCL_step2/.style={sci,precision=2,column name={$\Delta C_L$}},
 columns/dCD_step2/.style={sci,precision=2,column name={$\Delta C_D$}},
 columns/state_diff_step2_over_eps/.style={sci,precision=2,column name={$\|\delta\|_2/\|e\|_2$}},
 columns/state_diff_step2_over_eps_volume_weighted/.style={sci,precision=2,column name={$\|\delta\|_M/\|e\|_M$}}
]{\DataDir/naca_cancellation.csv}
\end{table}

The initial residual probes in this additional test use different trial states. For a trial displacement $q$, their difference has the form
\begin{equation}
 \Delta\Phi=\frac{3M_Bq}{2\Delta t}
       +F_B(Y+q)-F_B(Y)-\frac{M_B\eta^n}{\Delta t}
 \label{eq:unequal-trial-residual}
\end{equation}
With $q=e/4$ and $\eta^n=0$, the temporal part is three quarters of the older-only forcing magnitude. The measured ratio near 0.75 is therefore not a common-trial verification of Eq.~\eqref{eq:history-residual-identity}: the spatial residual difference also remains. Later cancellation is not expected even in the ideal first-step case. After $\delta^{n+1}=0$, the older level still contains $e/4$, giving a subsequent history combination $-e/8$.

\subsection{Perturbation-amplitude dependence}

Scaling the discrepancy by $\lambda=1,1/2,1/4$ reduces the Euclidean-state and lift superposition remainders for BAR2 and GP2 (Table~\ref{tab:naca-amplitude-app}). The Euclidean absolute state remainder decreases by factors of about 3.2--3.5 per halving, consistent with a higher-than-linear reduction over this finite range. The volume-weighted absolute remainder decreases much more slowly, from approximately $2.4\times10^{-5}$ to $1.5$--$1.6\times10^{-5}$, while its normalized value rises from about 3.4\% to 8.7--9.2\%. These results support directional response in selected diagnostics, but not a uniform quadratic remainder or a verified numerical floor in all norms.

\begin{table}[htbp]
\centering\footnotesize
\caption{Amplitude dependence of first-step superposition. $\mathcal E=\delta_f-\delta_n-\delta_{n-1}$ is the state remainder and $\mathcal E_L$ its lift counterpart. Relative remainders use the corresponding both-histories response. The Euclidean and volume-weighted state metrics have different amplitude dependence.}
\label{tab:naca-amplitude-app}
\pgfplotstabletypeset[publication table,
 columns={method,lambda,state_remainder_abs_euclid,state_superposition_defect_euclid,state_remainder_abs_M,state_superposition_defect_M,CL_remainder_abs},
 columns/method/.style={string type,column name={Transfer}},
 columns/lambda/.style={fixed,precision=2,column name={$\lambda$}},
 columns/state_remainder_abs_euclid/.style={sci,precision=2,column name={$\|\mathcal E\|_2$}},
 columns/state_superposition_defect_euclid/.style={sci,precision=2,column name={$\|\mathcal E\|_2/\|\delta_f\|_2$}},
 columns/state_remainder_abs_M/.style={sci,precision=2,column name={$\|\mathcal E\|_M$}},
 columns/state_superposition_defect_M/.style={sci,precision=2,column name={$\|\mathcal E\|_M/\|\delta_f\|_M$}},
 columns/CL_remainder_abs/.style={sci,precision=2,column name={$|\mathcal E_L|$}}
]{\DataDir/naca_amplitude_sequence.csv}
\end{table}
\FloatBarrier

\section{RANS numerical sensitivity}
\label{app:rans-controls}

\subsection{Inner iterations and time-step refinement}

Table~\ref{tab:rans-inner-app} distinguishes achieved residual criteria from iteration limits. The production rows cover the full repeated-remeshing histories; the other rows cover short controls. Raising the SA cap to 1600 reaches the density criterion throughout the control window. SA also reaches that criterion at every half- and quarter-step point. SST reaches its cap on 28--30 of 38 steps even with 6400 inner iterations, and on 57--58 of 78 half-step points. Increasing the iteration allowance is therefore not equivalent to obtaining a converged SST reference.

\begin{table}[htbp]
\centering\footnotesize
\caption{Inner-iteration behavior for the paired seed-1 calculations. Entries separated by a slash are GP2/NN2. $N_{\max}$ is the iteration allowance; residuals are median terminal $\log_{10}(\mathrm{rms}\,\rho)$. Production rows cover the full interval; the remaining rows cover the short numerical-control window.}
\label{tab:rans-inner-app}
\begin{tabular}{lrrrrr}
\toprule
Closure & $N_{\max}$ & $\Delta t/\Delta t_0$ & Physical steps & Steps at cap & Residual\\
\midrule
\InnerRow{SA}{production}{800}{1}
\InnerRow{SA}{cap1600}{1600}{1}
\InnerRow{SA}{halfstep}{800}{$1/2$}
\InnerRow{SA}{quarterstep}{800}{$1/4$}
\InnerRow{SST}{production}{800}{1}
\InnerRow{SST}{cap1600}{1600}{1}
\InnerRow{SST}{cap6400}{6400}{1}
\InnerRow{SST}{halfstep}{800}{$1/2$}
\bottomrule
\end{tabular}
\end{table}

Table~\ref{tab:rans-controls-app} reports the effect of increased caps using Eq.~\eqref{eq:control-sensitivity}. SA lift and total-drag responses change by at most 0.3\% of their own RMS; their smaller paired differences change by 8.4--10.9\%. For SST, the 6400-cap comparison changes the paired lift difference by 68.1\% of its production RMS. Persistent residual stagnation accompanies this sensitivity, but does not identify a unique cause in the turbulence equations.

\begin{table}[htbp]
\centering\small
\caption{Effect of increasing the inner-iteration cap, relative to the production setting. Columns have the same definition as Table~\ref{tab:rans-temporal-main}: individual values are the larger NN2 or GP2 sensitivity, and paired values refer to their difference. All entries are percentages.}
\label{tab:rans-controls-app}
\begin{tabular}{llrrrr}
\toprule
 & & \multicolumn{2}{c}{Individual histories} & \multicolumn{2}{c}{Paired difference}\\
\cmidrule(lr){3-4}\cmidrule(lr){5-6}
Closure & Iteration cap & $C_L$ & $C_D$ & $C_L$ & $C_D$\\
\midrule
\ControlRow{SA}{cap1600}{1600}
\ControlRow{SST}{cap1600}{1600}
\ControlRow{SST}{cap6400}{6400}
\bottomrule
\end{tabular}
\end{table}

For the SA step-refinement sequence, the paired lift RMS is approximately $1.21\times10^{-5}$, $1.60\times10^{-5}$, and $2.08\times10^{-5}$. The corresponding drag values are $2.19\times10^{-6}$, $3.77\times10^{-6}$, and $4.55\times10^{-6}$. The detailed paired response is not settled. Individual RMS amplitudes remain closer across the same samples, but that observation cannot replace a convergence test of the histories. Friction drag is more sensitive than total drag: the SA quarter-step history differences are approximately 76--91\% of the production friction-response RMS. No componentwise convergence claim follows from the smaller total-drag sensitivity.

\subsection{Steady mesh sensitivity and cross-mesh robustness}

The three steady mesh levels in Table~\ref{tab:rans-meshlevels-app} preserve the first off-wall spacing while changing the remaining resolution. Lift, total drag, and pressure drag decrease monotonically, but friction drag does not. The count-based spacing is only a mesh-density descriptor for this anisotropic family. The results do not demonstrate a common asymptotic error expansion, and no extrapolated loads or discretization-uncertainty estimate are reported. Differences in absolute loads remain distinct from errors in a paired transfer comparison on the same mesh.

\begin{table}[htbp]
\centering\small
\caption{Steady no-transfer loads on the three viscous meshes. $h/h_0=\sqrt{N_0/N}$ describes the change in total vertex count; it is not a uniform refinement ratio in every direction. The same first off-wall spacing is used on all levels.}
\label{tab:rans-meshlevels-app}
\pgfplotstabletypeset[publication table,
 columns={closure,level,nodes,h_ratio,CL,CD,CDp,CDv},
 columns/closure/.style={string type,column name={Closure}},
 columns/level/.style={string type,column name={Mesh}},
 columns/nodes/.style={fixed,precision=0,column name={Vertices}},
 columns/h_ratio/.style={fixed,precision=3,column name={$h/h_0$}},
 columns/CL/.style={fixed,precision=5,column name={$C_L$}},
 columns/CD/.style={fixed,precision=6,column name={$C_D$}},
 columns/CDp/.style={fixed,precision=6,column name={$C_{D,p}$}},
 columns/CDv/.style={fixed,precision=6,column name={$C_{D,v}$}}
]{\GeneratedDir/reassessed_rans_meshlevels.csv}
\end{table}

On the intermediate mesh, the short SA remeshing window gives individual lift RMS responses near $1.07\times10^{-4}$ and drag responses near $5.39\times10^{-5}$; the paired differences are $8.77\times10^{-6}$ and $9.94\times10^{-7}$, respectively. Thus the two methods also remain relatively close on another mesh. Because new node positions produce a different random displacement field, this comparison combines mesh resolution and disturbance realization. The smaller response cannot be attributed to refinement alone. Five to seven of its 38 physical steps also reach the production inner cap.

\begingroup
\small
\printbibliography
\endgroup
\end{document}
